\section{Aztec}
\label{sec:aztec}

Aztec es una red de capa 2 (\textit{Layer 2}) construida sobre Ethereum cuyo objetivo principal es
ofrecer contratos inteligentes con estado y ejecución privados~\cite{aztecdocs}. Su desarrollo
está a cargo de Aztec Labs, la misma organización que creó el lenguaje Noir y el sistema de
pruebas Barretenberg descritos en la Sección~\ref{sec:noir}. La red principal de Aztec, denominada
\textit{Ignition Chain}, comenzó a producir bloques sobre Ethereum en noviembre de
2025~\cite{theblock2025aztec}. El estudio de Aztec resulta pertinente para este trabajo por dos
razones. En primer lugar, constituye el entorno para el cual se concibieron las herramientas de
Noir, de modo que su arquitectura explica muchas de las decisiones de diseño de ese lenguaje. En
segundo lugar, Aztec es un sistema en el que la validez de las transiciones de estado de una
cadena se certifica ante otra mediante una prueba de conocimiento cero verificada por un contrato
inteligente, que es precisamente el principio sobre el que se construye el puente propuesto. La
presente sección describe el modelo de \textit{rollup} en el que se inscribe Aztec, su
representación del estado, la estructura de sus bloques, sus mecanismos de consenso y el
protocolo de mensajería que lo comunica con Ethereum.

\subsection{Rollups y pruebas de validez}
\label{sec:aztec_rollups}

Un \textbf{rollup} es una cadena que ejecuta transacciones fuera de una cadena principal, a la
que se denomina capa 1, y publica en ella tanto los datos necesarios para reconstruir su estado
como un compromiso sobre el estado resultante. La capa 1 actúa así como fuente de verdad y como
garante de la disponibilidad de los datos, mientras que el cómputo se realiza fuera de ella. Los
rollups se clasifican según el mecanismo con el que la capa 1 llega a aceptar un nuevo estado.

En los \textbf{rollups optimistas}, los nuevos estados se publican sin demostración y se presumen
válidos, salvo que algún participante presente una \textit{prueba de fraude} durante un período
de impugnación, que en las implementaciones habituales se extiende por siete días. Optimism,
cuya red de prueba Optimism Sepolia actúa como cadena de destino en este trabajo, pertenece a esta
categoría~\cite{optimismfaultproofs}. En los \textbf{rollups de validez}, también llamados
\textit{zk-rollups}, cada nuevo estado se acompaña de una prueba sucinta de que resulta de aplicar
correctamente las reglas del protocolo a un conjunto de transacciones, y un contrato verificador
desplegado en la capa 1 solo acepta el nuevo estado si la prueba es válida. Esta segunda
categoría elimina el período de impugnación y reemplaza la suposición de que algún participante
honesto vigila la cadena por una garantía criptográfica. Aztec pertenece a la categoría de
rollups de validez y añade a ella una propiedad que los demás rollups no ofrecen, que es la
posibilidad de ocultar a los propios operadores de la red el contenido de las transacciones
privadas~\cite{aztecdocs}.

\subsection{Estado privado y estado público}
\label{sec:aztec_estado}

A diferencia de Ethereum, cuyo estado es enteramente público, Aztec mantiene dos clases de estado
con modelos de datos distintos~\cite{aztecstate}. El \textbf{estado público} funciona de manera
análoga al de Ethereum, ya que cada contrato dispone de posiciones de almacenamiento cuyos valores
son visibles para todos los participantes y pueden ser modificados directamente por la ejecución
pública. El \textbf{estado privado}, en cambio, se representa mediante un modelo de salidas no
gastadas semejante al de Bitcoin, en el que la unidad de información es la \textit{nota}.

\begin{definition}
	\label{def:nota}
	Una \textbf{nota} es un registro cifrado que contiene una porción del estado privado de un
	contrato, como el saldo de un token perteneciente a un usuario. La red no almacena la nota,
	sino un compromiso criptográfico sobre ella, denominado \textbf{hash de nota}, que se calcula a
	partir de su contenido, su propietario, la posición de almacenamiento a la que pertenece, la
	dirección del contrato y un valor aleatorio. Solo el propietario de la nota conoce su
	contenido y puede descifrarla.
\end{definition}

Puesto que modificar un valor existente revelaría qué registro se actualizó, el estado privado
nunca se modifica ni se elimina. Para consumir una nota, su propietario publica un
\textbf{anulador} (\textit{nullifier}), que es un valor derivado de la nota y de una clave secreta
del propietario, calculado de forma que ningún observador pueda vincularlo con el hash de nota
correspondiente. Una transferencia privada consiste, en consecuencia, en anular las notas del
emisor y crear notas nuevas para el receptor y para el cambio, de manera similar al uso de dinero
en efectivo. La red rechaza toda transacción que intente publicar un anulador ya existente, lo que
impide gastar dos veces la misma nota sin necesidad de identificarla~\cite{aztecstate}.

Aztec adopta, además, la abstracción de cuentas de forma nativa. Toda cuenta es un contrato
inteligente que define sus propias reglas de autorización, de modo que el esquema de firma, la
gestión del \textit{nonce} y la forma de pago de las tarifas quedan a criterio de quien
programa la cuenta~\cite{aztecdocs}. Por ejemplo, una cuenta puede verificar firmas ECDSA sobre la
curva secp256k1, que es la misma que emplea Ethereum, u otro esquema distinto.

\subsection{Árboles de estado}
\label{sec:aztec_arboles}

El estado de Aztec se organiza en cinco árboles de Merkle binarios, cuyas raíces conforman el
compromiso sobre el estado global de la red. A diferencia de Ethereum, que anida un árbol de
almacenamiento dentro de cada cuenta (Sección~\ref{sec:eth_mpt}), Aztec almacena los valores
directamente como hojas de árboles globales~\cite{aztecstate}. Los árboles emplean funciones de
hash algebraicas, como Poseidon2, que se expresan como circuito con un número de restricciones muy
inferior al de Keccak-256 o SHA-256. La Tabla~\ref{tab:aztec_arboles} resume su función y su
altura, según los parámetros vigentes del protocolo~\cite{aztecconstants}.

\begin{table}[htbp]
	\centering
	\caption{Árboles de estado de Aztec.}
	\label{tab:aztec_arboles}
	\small
	\begin{tabular}{>{\raggedright\arraybackslash}p{3.8cm} >{\raggedright\arraybackslash}p{8.6cm} >{\centering\arraybackslash}p{1.4cm}}
		\toprule
		\textbf{Árbol} & \textbf{Contenido} & \textbf{Altura} \\
		\midrule
		Árbol de hashes de nota & Compromisos sobre las notas del estado privado. Solo admite
		inserciones al final. & 42 \\
		\addlinespace
		Árbol de anuladores & Anuladores publicados, organizados como árbol de Merkle indexado. & 42 \\
		\addlinespace
		Árbol de datos públicos & Valores del almacenamiento público de todos los contratos. & 40 \\
		\addlinespace
		Árbol de mensajes de L1 a L2 & Mensajes enviados desde Ethereum hacia Aztec. & 36 \\
		\addlinespace
		Árbol de archivo & Hashes de los encabezados de todos los bloques anteriores. & 30 \\
		\bottomrule
	\end{tabular}
\end{table}

El árbol de anuladores plantea un requisito particular. Para insertar un anulador, la red debe
demostrar que este no existe previamente en el árbol, lo cual exige una \textit{prueba de no
pertenencia}. Un árbol de Merkle disperso, en el que cada valor posible ocupa la hoja de índice
igual a dicho valor, permitiría esta prueba de forma directa, pero requeriría una profundidad de
254 niveles para cubrir los elementos del cuerpo de la curva BN254, con el consiguiente costo en
restricciones. Aztec resuelve este problema con un \textbf{árbol de Merkle
indexado}~\cite{aztecindexedtree}, en el que cada hoja almacena una terna
$(v, i_{\mathrm{sig}}, v_{\mathrm{sig}})$ formada por el valor $v$, el índice $i_{\mathrm{sig}}$ de
la hoja que contiene el siguiente valor mayor y dicho valor $v_{\mathrm{sig}}$. Las hojas forman,
así, una lista enlazada ordenada. Para demostrar que un valor $x$ no pertenece al árbol basta con
exhibir la pertenencia de la hoja que lo precede, denominada \textit{anulador inferior}, y
comprobar que $v < x < v_{\mathrm{sig}}$, o bien que $v < x$ y $v_{\mathrm{sig}} = 0$ cuando $x$
es mayor que todos los valores almacenados. Como la posición de una hoja ya no depende de su
valor, una profundidad moderada es suficiente.

El \textbf{árbol de archivo} cumple una función semejante a la del encadenamiento por hash
descrito en la Sección~\ref{sec:arquitectura_blockchain}. Al almacenar el hash de cada
encabezado, permite demostrar dentro de un circuito, con una prueba de inclusión logarítmica, que
un bloque anterior forma parte de la historia de la cadena.

\subsection{Ejecución de transacciones y estructura de bloques}
\label{sec:aztec_bloques}

La ejecución de una transacción en Aztec se divide en una fase privada y una fase
pública~\cite{aztectransactions}. La fase privada se ejecuta en el dispositivo del usuario,
dentro de un componente denominado PXE (\textit{Private eXecution Environment}), que almacena las
claves y las notas del usuario, ejecuta las funciones privadas de los contratos y genera las
pruebas correspondientes. Cada llamada a una función privada es procesada por el \textbf{circuito
de núcleo privado} (\textit{private kernel}), que verifica la prueba de la función, comprueba el
contexto de la llamada y acumula sus efectos, es decir, los hashes de nota, los anuladores, los
mensajes hacia Ethereum y las llamadas públicas pendientes~\cite{aztecprivatekernel}. Este
circuito se aplica de forma recursiva, una vez por llamada, y concluye con una prueba única que
certifica toda la ejecución privada de la transacción. Es el único circuito del protocolo que
requiere efectivamente la propiedad de conocimiento cero, puesto que debe ocultar qué contrato y
qué función se ejecutaron, quién es el usuario y cuáles fueron las entradas y salidas. Los
restantes circuitos utilizan pruebas sucintas para comprimir la verificación, pero no necesitan
ocultar sus testigos.

El usuario envía a la red la prueba de la fase privada junto con sus efectos. La fase pública,
que comprende las funciones públicas encoladas por la fase privada, se ejecuta posteriormente en
los nodos de la red mediante la AVM (\textit{Aztec Virtual Machine}), una máquina virtual
conceptualmente análoga a la EVM, cuya ejecución también se demuestra mediante una prueba. El
costo de las transacciones se mide en \textit{mana}, unidad equivalente al gas de Ethereum
(Sección~\ref{sec:eth_gas}), que se desglosa en un componente de ejecución y otro de
disponibilidad de datos, y cuyo precio se determina mediante un mecanismo inspirado en
EIP-1559~\cite{aztecfees}.

El encabezado de un bloque de Aztec contiene la raíz del árbol de archivo antes del bloque, las
raíces de los árboles de hashes de nota, de anuladores, de datos públicos y de mensajes de L1 a
L2 tras su ejecución, un conjunto de variables globales que incluye el identificador de la
cadena, el número de bloque, el número de \textit{slot}, la marca de tiempo y los precios del
mana, y los totales de tarifas y de mana consumidos~\cite{aztecconstants}. Los bloques se agrupan
en \textit{checkpoints}, que son las unidades que se publican en Ethereum, y estos a su vez en
\textit{épocas} de hasta 32 checkpoints.

La validez de una época completa se demuestra mediante una única prueba, construida por los
\textbf{circuitos de rollup} a partir de un árbol binario de pruebas~\cite{aztecrollupcircuits}.
En el nivel inferior, un circuito base verifica la prueba de cada transacción y actualiza los
árboles de estado. Los niveles superiores fusionan recursivamente pares de pruebas, primero en
pruebas de bloque, luego de checkpoint y finalmente en la prueba raíz de la época, cuyas entradas
públicas incluyen las raíces del árbol de archivo anterior y posterior. Esta topología permite
distribuir la generación de las pruebas entre numerosos probadores que trabajan en paralelo. La
prueba raíz se verifica en el contrato de rollup desplegado en Ethereum, de modo que el costo de
verificación en la capa 1 es independiente del número de transacciones incluidas en la época. La
Figura~\ref{fig:aztec_ciclo} resume el recorrido de una transacción desde el dispositivo del
usuario hasta su asentamiento en Ethereum.

\begin{figure}[htbp]
	\centering
	\begin{tikzpicture}[
		node distance=0.55cm,
		caja/.style={draw, rounded corners=2pt, align=center, minimum width=8.2cm,
			minimum height=0.95cm, font=\small},
		flecha/.style={-{Stealth[length=2.2mm]}, thick},
		etiqueta/.style={font=\footnotesize\itshape, align=left, anchor=west}
		]
		\node[caja] (pxe) {PXE del usuario\\
			(ejecución privada y núcleo privado)};
		\node[caja, below=of pxe] (secuenciador) {Secuenciador y comité\\
			(ejecución pública en la AVM y atestación del bloque)};
		\node[caja, below=of secuenciador] (l1datos) {Contrato de rollup en Ethereum\\
			(publicación del checkpoint y sus datos)};
		\node[caja, below=of l1datos] (probadores) {Probadores\\
			(circuitos de rollup y prueba raíz de la época)};
		\node[caja, below=of probadores] (l1prueba) {Contrato de rollup en Ethereum\\
			(verificación de la prueba y asentamiento)};

		\draw[flecha] (pxe) -- (secuenciador)
			node[etiqueta, midway, xshift=4.25cm] {prueba privada y efectos};
		\draw[flecha] (secuenciador) -- (l1datos)
			node[etiqueta, midway, xshift=4.25cm] {checkpoint atestiguado};
		\draw[flecha] (l1datos) -- (probadores)
			node[etiqueta, midway, xshift=4.25cm] {época concluida};
		\draw[flecha] (probadores) -- (l1prueba)
			node[etiqueta, midway, xshift=4.25cm] {prueba de validez};
	\end{tikzpicture}
	\caption{Recorrido de una transacción de Aztec desde su ejecución en el dispositivo del
	usuario hasta su asentamiento en Ethereum.}
	\label{fig:aztec_ciclo}
\end{figure}

\subsection{Secuenciadores, probadores y finalidad}
\label{sec:aztec_consenso}

La producción de bloques en Aztec está a cargo de \textbf{secuenciadores}, que son participantes
que depositan una garantía de 200\,000 tokens AZTEC para integrar el conjunto de validadores de la
red~\cite{aztecoperators}. En cada slot, se selecciona de forma pseudoaleatoria un proponente a
partir de la aleatoriedad provista por Ethereum, y un comité de secuenciadores reejecuta las
transacciones propuestas y firma atestaciones si las considera válidas~\cite{aztecblocks}. Un
bloque solo se acepta si reúne las atestaciones de al menos dos tercios más uno de los miembros
del comité, y la publicación de un bloque sin ese respaldo constituye una falta sancionable con
la pérdida de parte de la garantía.

Las atestaciones proporcionan una confirmación provisional, pero no constituyen la garantía
definitiva. Una vez concluida una época, los \textbf{probadores} generan la prueba de validez de
todos sus bloques y la remiten al contrato de rollup, que la verifica y actualiza el estado
asentado. Si ninguna prueba válida llega dentro del plazo previsto, los bloques no demostrados se
descartan y sus transacciones deben volver a proponerse~\cite{aztecblocks}. La seguridad de las
transiciones de estado asentadas en Ethereum descansa, por lo tanto, en la validez de la prueba y
no en la honestidad del comité, cuyo papel se limita a ordenar las transacciones y a garantizar la
disponibilidad de los datos.

\subsection{Mensajería entre Ethereum y Aztec}
\label{sec:aztec_mensajeria}

Las aplicaciones de Aztec se comunican con Ethereum mediante un protocolo de mensajes asíncronos.
El extremo de una aplicación en Ethereum se denomina \textbf{portal}, y es un contrato inteligente
vinculado a una dirección específica de Aztec~\cite{aztecmessaging}. Los mensajes desde Ethereum
hacia Aztec se depositan en un contrato denominado \textit{Inbox}, que almacena únicamente un
compromiso de cada mensaje. Al procesar el checkpoint siguiente, los circuitos de rollup insertan
esos mensajes en el árbol de mensajes de L1 a L2, y un contrato de Aztec puede consumirlos
demostrando su pertenencia a dicho árbol y publicando un anulador que impide consumirlos dos
veces. En sentido inverso, los mensajes que emite una transacción de Aztec quedan disponibles en
un contrato denominado \textit{Outbox} una vez que la prueba de la época que los contiene ha sido
verificada en Ethereum, momento a partir del cual el portal correspondiente puede consumirlos.

Este protocolo permite construir puentes de activos del tipo \textit{lock-and-mint} sin
intermediarios de confianza. El propio protocolo incluye uno de ellos, el contrato
\texttt{FeeJuicePortal}, que bloquea en Ethereum los tokens destinados al pago de tarifas y envía
un mensaje al contrato correspondiente de Aztec, que acredita el saldo al destinatario cuando este
presenta el secreto asociado al depósito~\cite{aztecfees}. La acuñación en Aztec no depende de un
\textit{relayer}, sino de que el mensaje figure en un árbol cuya raíz ha sido establecida por los
circuitos de rollup y verificada en Ethereum, y la liberación de fondos en Ethereum depende, a su
vez, de una prueba de validez aceptada por el contrato de rollup.

Esta garantía es posible porque Aztec y Ethereum comparten un contrato de asentamiento, ya que es
Ethereum quien verifica las pruebas de Aztec y, por ende, conoce con certeza su estado. Dos
cadenas independientes, como Ethereum Sepolia y Optimism Sepolia, no disponen de un mecanismo
equivalente, puesto que ninguna verifica de forma nativa el consenso de la otra. El puente
desarrollado en este trabajo traslada el principio de los rollups de validez a este escenario,
reemplazando la verificación de la prueba de un rollup por la verificación de una prueba de
conocimiento cero sobre la firma del comité de sincronización de Ethereum y sobre la inclusión del
evento de bloqueo, según la cadena de compromisos descrita en la
Sección~\ref{sec:eth_cadena_verificacion}.
